\chapter{8}



Kurz vor Sonnenaufgang schrak Louis hoch. 
Verwirrt blinzelte er sich von Ast zu Ast nach oben. Hatte er geschlafen? Die ersten
Lastwagen des Tages schossen nicht weit entfernt über die Landstrasse.
Und blitzschnell war die Erinnerung zurück. Die Bilder kamen, einige
messerscharf, ausser - er suchte danach, strengte sich an. Was war in
ihn gefahren? Was war dort oben passiert? Warum hatte sie ihn provoziert?
Da waren sie wieder, die Fragen ohne Antworten.
Louis stand entschlossen auf.
Sein Atem wurde schneller. Und wieder flammte eine Spur Zorn auf.
Mit kurzem Blick zurück sah er den kleinen Stock im Laub liegen, bückte
sich, hob ihn auf und steckte ihn mit befremdender Selbstverständlichkeit in den Rucksack.
Er würde es als Anhalter versuchen, sich an die Autobahnauffahrt bei der
naheliegenden Raststätte stellen und warten. Die Entscheidung kam
überzeugend.

\sectionbreak

Louis suchte die Waschräume der Raststätte auf. Die Herrentoilette war
leer.
Gedankenversunken wusch er sich die Hände und begann seine Trinkflasche
mit Wasser zu füllen. Bis er in den Spiegel schaute und zusammenfuhr.
Wie er nur aussah, schmutzig, heruntergekommen. Der Anblick machte ihm
Angst.
Er begann sich Wasser ins Gesicht zu spritzen, spülte die Zähne, rückte
sein T-Shirt zurecht und fuhr sich mit den nassen Händen immer wieder
hektisch durchs Haar. Er brauchte Haargelée, vielleicht eine Brille und vor allem 
andere Kleidung.

Louis betrat den Shop. Der war gut besucht. Bemüht lässig bewegte er sich
zwischen den Regalen. Nach dem Einkauf verschwand er zum zweiten Mal in die Herrentoilette. 
Er kreuzte einen älteren Mann, der eben den Raum verliess,
huschte an ihm vorbei und stellte sich an einen der Waschtische. Er war
alleine. Mit raschen Bewegungen strich er sich das Gelée ins Haar und
kämmte es streng nach hinten. Dann schloss er sich in einer Kabine ein.
Er zog sich das neue «Milano-Bella Città»-Shirt über und setzte die
Brille auf. Lesebrille 1.0. Die einzige Möglichkeit. Besser als nichts.
Beim Blick durch die Gläser starrte er die Wände entlang in den
verzerrten Raum. Er versuchte über den Rand zu blicken und setzte sich
die Baseball-Kappe auf. Schliesslich zog er die Laufschuhe aus und schlüpfte in die neuen
Strandschlappen. Die abgelegten Sachen stopfte er in den Rucksack. Beim
Blick in den Spiegel staunte er. Am stärksten schien ihn die Brille zu
verändern. Er sah aus wie ein läppischer Tourist. Es musste an 
diesem T-Shirt und der einfältigen Mütze liegen. Immerhin, die Verwandlung schien
gelungen. Louis zwang sich ein blasses Lächeln ab. Jetzt bestimmte er. Niemand
anderes als er. Augenblicklich bewegten sich seine Finger. Er sehnte sich nach seiner Gitarre und wünschte nichts mehr, 
als zu spielen. Gerade jetzt und ganz für sich.

\sectionbreak

Auf dem kurzen Weg zur Autobahnauffahrt schlängelte er sich auf dem
weiträumigen, lärmigen Parkplatz zwischen Autos und Lastwagen hindurch.
Plötzlich fand er sich vor der offenen Tür eines massigen Gefährts.
